\section{対象のデータと用語}
\label{sec_data}
本研究で用いるデータベースは，ボケてから収集したお題画像6,823,270件と，それらに投稿されたボケ約9,250万件からなる．
ボケには，本文のほかに，お題の ID，投稿者，評価の合計（星の数）などの情報が含まれる．
サイトには不適切な投稿を示す \texttt{unsafe} と笑いの分類を示す \texttt{category} の項目もあるが，予備調査では \texttt{unsafe} はボケの0.03\%にしか付いておらず，\texttt{category} は89\%が既定値であったため，ラベルとしては用いなかった．
また，評価が5件以上付いたボケは全体の2〜4\%にとどまり，評価のないボケの多くは面白くないのではなく閲覧されていないと考えられる．

各お題画像には，筆者らの先行の作業により，大喜利の形式，画像中の文字の有無，画像の媒体，写真の簡潔な説明文（キャプション）が付与されている．
大喜利の形式は次の4つである．
\begin{itemize}
    \item 通常：画像に自由に一言を付ける．
    \item カルタ：画像の隅に書かれたひらがな1文字（札のかな）から始まる一言を答える．
    \item 穴埋め：画像中の文の空欄に入る語句を答える．
    \item お題文：画像に書かれた問いや指示に答える．
\end{itemize}
画像の媒体は，写真，アニメ，ゲーム，漫画，文字主体の画像，絵画，イラストの7つであり，複数を同時に持ちうる．
キャプションは写真にのみ付与されており，人物の名前は書かない設計である．
本稿では，これらをお題のラベルと呼び，ボケのラベルを推定する際の入力として用いる．

ボケの長さは形式によって異なる．
形式と媒体で分けた7つの層（\ref{sec_label}章）から300件ずつ抽出したお題の全てのボケ（約2.9万件）では，文字数の中央値は穴埋めが8，お題文が10，カルタが12，通常が14〜16であった．
穴埋めのボケは空欄に入る語だけのことが多く，お題画像の文と合わせなければ意味が分からない．

本稿では，Qwen3-VL-32B のように大規模な VLM を教師，教師が付けたラベル（教師ラベル）で学習する高速な分類器を生徒と呼ぶ．
目視ラベルは，画像とボケを表示して内容を確認できる AI エージェントである Claude Code が，その時点の定義書に従って付与した．
付与の際は，お題画像とボケを8組程度ずつ並べた一覧画像を作成し，評価する手法の予測は表示しなかった．
目視ラベルは全て同一の判断者によるものであり，複数の判断者の間の一致度は測っていない．
この点は本研究の限界であり，\ref{sec_discussion}章で改めて述べる．
